% ── CHAPTER 2: FROM ARISTOTLE TO WHITEHEAD: WHY SUBSTANCE ONTOLOGY FAILS (~8,000 words) ──
\chapter{From Aristotle to Whitehead: Why Substance Ontology Fails}
\label{ch:aristotle-whitehead}

% STATUS: TODO
% TARGET: ~8,000 words
% PRIMARY SOURCES:
%   pub2/sections/02-against-substance.tex  — substance ontology argument (academic)
%   pub2/sections/03-whitehead-process.tex  — Whitehead's four categories (academic)
% SECONDARY SOURCES: Aristotle_Metaphysics, Whitehead1929, Whitehead1933,
%                    Rescher1996, Griffin2000, Epperson2004, Cobb1965
% HARD RULES: See BOOK_SHELL.md
% REGISTER: More philosophical than Ch. 1 but still accessible.
%           Introduce technical terms (actual occasion, prehension) with analogies.

% \epigraph{It is more important that a proposition be interesting than
%   that it be true.}{Alfred North Whitehead, \textit{Process and Reality}}

% ─────────────────────────────────────────────────────────────────────────────
%  OPENING (~500 words)
% ─────────────────────────────────────────────────────────────────────────────
% TODO: Chapter 1 showed what physics discovered. This chapter asks: what kind
% of philosophy can accommodate it? Spoiler: not the one we inherited.
% Set up the central move: substance ontology is the default Western metaphysics;
% quantum mechanics is formally incompatible with it; Whitehead's process
% ontology is the most fully developed alternative.

% ─────────────────────────────────────────────────────────────────────────────
\section{The Substance Tradition: Aristotle to Newton}
\label{sec:ch2-substance}
% (~1,500 words)
% ─────────────────────────────────────────────────────────────────────────────
% TODO: History of substance ontology for a general reader.
% KEY POINTS:
%   - Aristotle: substances are primary; properties and relations are secondary
%   - A substance is what persists through change; its properties are predicated
%     of it but it itself is not predicated of anything else
%   - Newton: the billiard ball as the perfect substance — mass, position,
%     trajectory all inherent in the ball
%   - Descartes: mind and body as two different kinds of substance
%   - The deep assumption: THINGS first; relations and processes second
% SOURCE: pub2 §2.1 (substance ontology), Aristotle_Metaphysics, Whitehead1929 §1
%         Translate academic prose to accessible narrative.
% CITE: Aristotle_Metaphysics

% ─────────────────────────────────────────────────────────────────────────────
\section{Why Substance Ontology Cannot Survive Bell's Theorem}
\label{sec:ch2-bell-substance}
% (~1,500 words)
% ─────────────────────────────────────────────────────────────────────────────
% TODO: Connect Ch. 1's physics finding to philosophical consequence.
% KEY POINTS:
%   - Substance ontology requires intrinsic properties (possessed independently
%     of relations to other things)
%   - Bell's theorem shows there are no such local intrinsic properties
%   - Therefore substance ontology is formally incompatible with quantum mechanics
%   - This is not a philosophical preference — it is a mathematical result
%   - Contrast with Bohm: Bohm restores determinism but at the cost of
%     fundamental nonlocality — the substance has moved to the wave function,
%     which is global, not local
% SOURCE: pub2 §2.2–2.4 (steps 2–4 of the argument); translate from academic.
% CITE: Bell1964, Aspect1982, Bohm1980

% ─────────────────────────────────────────────────────────────────────────────
\section{Whitehead's Inversion: Events Before Things}
\label{sec:ch2-whitehead-inversion}
% (~2,000 words)
% ─────────────────────────────────────────────────────────────────────────────
% TODO: Introduce process ontology and its central inversion.
% KEY POINTS:
%   - Whitehead's starting point: what we perceive as "things" are stable
%     patterns of relational events — not substances that undergo change
%   - The inversion: EVENTS are primary; things are derivative
%   - Actual occasion: the fundamental unit of reality — a momentary event of
%     becoming, not an enduring substance
%   - "The event of becoming is more fundamental than the thing that becomes"
%   - The stone is not a thing; it is a high-frequency pattern of relational
%     events occurring with sufficient regularity to appear stable
% ANALOGY: A flame — continuous process, not a substance. What persists is the
%          pattern of burning, not a burning-stuff.
% SOURCE: pub2 §3.1 (primacy of process), §3.2 (key categories); translate.
% CITE: Whitehead1929, Rescher1996

% ─────────────────────────────────────────────────────────────────────────────
\section{Prehension, Concresence, and Eros}
\label{sec:ch2-prehension}
% (~2,000 words)
% ─────────────────────────────────────────────────────────────────────────────
% TODO: Whitehead's four key concepts in accessible terms.
% KEY POINTS:
%   ACTUAL OCCASION: The momentary event of becoming. Arises, achieves its
%     definite character, perishes — contributing to the objective past.
%   PREHENSION: The relational act by which an occasion "grasps" prior occasions.
%     The way the past becomes present. Not perception — more fundamental.
%   CONCRESCENCE: The process by which an occasion integrates its prehensions
%     into a unified, definite character. The physics of self-constitution.
%   EROS: Whitehead's term for the creative drive toward richer, more complex
%     relational integration. The force that draws occasions toward higher
%     forms of unity. This is Whitehead's name for what the Bahá'í writings
%     call mahabbat.
% ANALOGY: Prehension = the way a conversation incorporates everything that's
%          been said. Concrescence = arriving at a response.
% SOURCE: pub2 §3.2–3.3 (key categories, eros); translate from academic.
% CITE: Whitehead1929, Whitehead1933, Griffin2000

% ─────────────────────────────────────────────────────────────────────────────
\section{Why This Is Not New-Age Philosophy}
\label{sec:ch2-not-newage}
% (~700 words)
% ─────────────────────────────────────────────────────────────────────────────
% TODO: Pre-empt the dismissal. Process philosophy is rigorous, has a century
% of development, and is endorsed by serious physicists and philosophers.
% KEY POINTS:
%   - Whitehead was a professional mathematician and logician
%   - Process and Reality is one of the most technically demanding works in
%     twentieth-century philosophy
%   - David Bohm, Henry Stapp, Michael Epperson, and others have developed
%     the connection to physics
%   - It is not a metaphor; it is a formal ontological framework
% SOURCE: Write fresh; Epperson2004, Stapp2007, Griffin2000 for citations.
% CITE: Epperson2004, Stapp2007

% ─────────────────────────────────────────────────────────────────────────────
%  CLOSING (~300 words)
% ─────────────────────────────────────────────────────────────────────────────
% TODO: Transition. Physics broke substance ontology; process ontology is the
% most rigorous replacement available. Next chapter: the Bahá'í writings arrived
% at the same framework a generation before Whitehead named it.

\begin{convergencemoment}{From Aristotle to Whitehead: Why Substance Ontology Fails}
% TODO: Fill in all three traditions' perspectives.
% From the quantum side: Bell's theorem formally closes substance ontology.
% From the process-philosophical side: Whitehead's process ontology is the
%   most fully developed alternative — 100 years of development.
% From the Bahá'í side: [bridge to Ch. 3 — will be stated there]
% One sentence statement: [write when drafting]
\end{convergencemoment}
